% Автоматаар үүсгэсэн — эх файл: lab01/README.md
\chapter{Лаб 1 — Хоёр давхаргат платформ ба ирмэг–үүлний гүүр}

\begingroup\scriptsize\begin{longtable}[]{@{}
  >{\raggedright\arraybackslash}p{(\columnwidth - 2\tabcolsep) * \real{0.5000}}
  >{\raggedright\arraybackslash}p{(\columnwidth - 2\tabcolsep) * \real{0.5000}}@{}}
\toprule
\endhead
\bottomrule
\endlastfoot
\textbf{7 хоног} & III \\
\textbf{Хугацаа} & 4 цаг \\
\textbf{Суралцахуйн үр дүн} & ҮД1 (үүлд суурилсан платформ зохиох), ҮД2 (өргөтгөх боломжтой backend) \\
\textbf{Үнэлгээ} & Бичгийн тайлан, 10 оноо \\
\textbf{Гол хэмжилт} & Гүүрний RTT, Pi-гийн үлдэгдэл санах ой, үйлчилгээний эхлэх хугацаа \\
\end{longtable}\endgroup

\hypertarget{ux437ux43eux440ux438ux43bux433ux43e}{%
\section{1. Зорилго}\label{ux437ux43eux440ux438ux43bux433ux43e}}

Энэ лабораторийн эцэст та \textbf{хоёр машин дээр тархсан, ажиллагаатай IoT платформ}-той болно. Гол нь стекийг асаах биш --- \textbf{хаана юу ажиллах ёстойг тоогоор нотлох} явдал.

Гурван асуултад тоон хариулт өгнө:

\begin{enumerate}
\def\labelenumi{\arabic{enumi}.}
\tightlist
\item
  Ирмэгийн давхарга Raspberry Pi 3B дээр \textbf{хэдий хэмжээний нөөц} зарцуулах вэ, хэдий хэр илүүдэлтэй вэ?
\item
  Ирмэг ба үүлний хооронд мессеж явахад \textbf{хэдэн миллисекунд} зарцуулагдах вэ?
\item
  Яагаад платформын давхаргыг (EMQX, InfluxDB, Grafana) Pi 3B дээр ажиллуулах \textbf{боломжгүй} вэ --- таамаглал биш, тооцоо?
\end{enumerate}

\hypertarget{ux443ux440ux44cux434ux447ux438ux43bux441ux430ux43d-ux43dux4e9ux445ux446ux4e9ux43b}{%
\section{2. Урьдчилсан нөхцөл}\label{ux443ux440ux44cux434ux447ux438ux43bux441ux430ux43d-ux43dux4e9ux445ux446ux4e9ux43b}}

\texttt{SETUP.md} бүрэн дуусгасан байх ёстой. Ялангуяа:

\begin{itemize}
\tightlist
\item
  {[}Pi{]} \texttt{uname\ -m} → \texttt{aarch64}
\item
  {[}Pi{]} \texttt{free\ -h} → Swap 2.0Gi, Mem \textasciitilde925Mi
\item
  {[}Pi{]} \texttt{vcgencmd\ get\_throttled} → \texttt{0x0}
\item
  {[}ЗК{]} Docker Desktop-д ≥ 6 GB
\item
  {[}ЗК{]} LAN IP мэдэгдэж байгаа
\end{itemize}

\begin{Shaded}
\begin{Highlighting}[]
\CommentTok{\# [Pi] Pi дээр}
\BuiltInTok{cd}\NormalTok{ \textasciitilde{}/cnc302/edge }\KeywordTok{\&\&} \FunctionTok{make}\NormalTok{ check      }\CommentTok{\# бүгд OK байх ёстой}
\end{Highlighting}
\end{Shaded}

Аль нэг ✘ байвал \textbf{энд зогсоод SETUP.md-д буц}. Лабораторийн 4 цагийг орчин тохируулахад зарцуулж болохгүй.

\hypertarget{ux43eux43dux43eux43bux44bux43d-ux441ux430ux43dux443ux443ux43bux433ux430}{%
\section{3. Онолын сануулга}\label{ux43eux43dux43eux43bux44bux43d-ux441ux430ux43dux443ux443ux43bux433ux430}}

\hypertarget{edgefogcloud-ux43aux43eux43dux442ux438ux43dux443ux443ux43c}{%
\subsection{Edge--Fog--Cloud континуум}\label{edgefogcloud-ux43aux43eux43dux442ux438ux43dux443ux443ux43c}}

IoT лавлах архитектурт өгөгдөл нэг чиглэлд урсдаггүй. Гурван шийдвэрийн цэг байна:

\begingroup\scriptsize\begin{longtable}[]{@{}lll@{}}
\toprule
Асуулт & Ирмэг дээр & Үүл дээр \\
\midrule
\endhead
\bottomrule
\endlastfoot
Хэдэн мс дотор хариу хэрэгтэй вэ & \textless{} 100 мс & \textgreater{} 1 сек болно \\
Холбоос тасарвал ажиллах ёстой юу & тийм & үгүй \\
Хэдий хэмжээний өгөгдөл боловсруулах вэ & багахан & асар их \\
Түүх хадгалах уу & үгүй (эсвэл түр) & тийм \\
\end{longtable}\endgroup

Энэ курсын байрлуулалт яг үүнийг дагана: \textbf{хурдан, локал, тасалдалд тэсвэртэй зүйл Pi дээр; хүнд, түүхэн, харьцуулсан зүйл компьютер дээр.}

\hypertarget{ux433ux4afux4afux440-bridge-ux433ux44dux436-ux44eux443-ux432ux44d}{%
\subsection{Гүүр (bridge) гэж юу вэ}\label{ux433ux4afux4afux440-bridge-ux433ux44dux436-ux44eux443-ux432ux44d}}

MQTT гүүр бол \textbf{брокер хоорондын клиент холболт}. Ирмэгийн mosquitto нь үүлний EMQX-д ердийн MQTT клиент шиг холбогдож, тодорхой сэдвийн мессежийг дамжуулна.

\begin{verbatim}
төхөөрөмж ──► mosquitto (Pi) ══гүүр══► EMQX (компьютер) ──► InfluxDB
                    │
                    └── диск: холбоос тасрахад мессежийг ХАДГАЛНА
\end{verbatim}

Гүүрний гурван чухал тохиргоо (\texttt{edge/\allowbreak{}mosquitto/\allowbreak{}conf.\allowbreak{}d/\allowbreak{}bridge.\allowbreak{}conf}):

\begingroup\scriptsize\begin{longtable}[]{@{}
  >{\raggedright\arraybackslash}p{(\columnwidth - 4\tabcolsep) * \real{0.3333}}
  >{\raggedright\arraybackslash}p{(\columnwidth - 4\tabcolsep) * \real{0.3333}}
  >{\raggedright\arraybackslash}p{(\columnwidth - 4\tabcolsep) * \real{0.3333}}@{}}
\toprule
\begin{minipage}[b]{\linewidth}\raggedright
Тохиргоо
\end{minipage} & \begin{minipage}[b]{\linewidth}\raggedright
Утга
\end{minipage} & \begin{minipage}[b]{\linewidth}\raggedright
Яагаад
\end{minipage} \\
\midrule
\endhead
\bottomrule
\endlastfoot
\texttt{topic\ cnc302/\allowbreak{}shutis/\allowbreak{}\#\ out\ 1} & юуг, хаашаа, QoS хэдээр & угтвар зөрвөл мессеж \textbf{хэзээ ч} үүлэнд хүрэхгүй \\
\texttt{cleansession\ false} & session хадгална & тасрахад дараалал үлдэнэ \\
\texttt{restart\_timeout\ 5\ 60} & дахин холбогдох завсарлага & 5→60 сек өсөх backoff \\
\end{longtable}\endgroup

\hypertarget{ux441ux430ux43dux430ux445-ux43eux439ux43d-ux445ux44fux437ux433ux430ux430ux440-ux431ux43eux43b-ux430ux440ux445ux438ux442ux435ux43aux442ux443ux440ux44bux43d-ux445ux44fux437ux433ux430ux430ux440}{%
\subsection{Санах ойн хязгаар бол архитектурын хязгаар}\label{ux441ux430ux43dux430ux445-ux43eux439ux43d-ux445ux44fux437ux433ux430ux430ux440-ux431ux43eux43b-ux430ux440ux445ux438ux442ux435ux43aux442ux443ux440ux44bux43d-ux445ux44fux437ux433ux430ux430ux440}}

ThingsBoard-ын JVM ганцаараа \texttt{-Xmx1500m} шаарддаг. Pi 3B-д \textbf{нийт} 925 MiB. Энэ бол тохиргооны асуудал биш --- \textbf{физик хязгаар}. Тиймээс платформын давхарга компьютер дээр байх нь тохиромж биш, зайлшгүй шаардлага.

\hypertarget{ux430ux43bux445ux43cux443ux443ux434}{%
\section{4. Алхмууд}\label{ux430ux43bux445ux43cux443ux443ux434}}

\hypertarget{ux430ux43bux445ux430ux43c-1-ux438ux440ux43cux44dux433ux438ux439ux43d-ux441ux443ux443ux440ux44c-ux445ux44dux43cux436ux438ux43bux442-25-ux43cux438ux43d-pi}{%
\subsection{Алхам 1 --- Ирмэгийн суурь хэмжилт (25 мин) {[}Pi{]}}\label{ux430ux43bux445ux430ux43c-1-ux438ux440ux43cux44dux433ux438ux439ux43d-ux441ux443ux443ux440ux44c-ux445ux44dux43cux436ux438ux43bux442-25-ux43cux438ux43d-pi}}

Юу ч асаахаас өмнө \textbf{хоосон Pi-гийн төлөвийг} бичиж авна. Дараагийн бүх тоо үүнтэй харьцуулагдана.

\begin{Shaded}
\begin{Highlighting}[]
\BuiltInTok{cd}\NormalTok{ \textasciitilde{}/cnc302}
\ExtensionTok{pkill} \AttributeTok{{-}f}\NormalTok{ vscode{-}server           }\CommentTok{\# VS Code сервер 150–250 MiB иднэ!}
\FunctionTok{sleep}\NormalTok{ 3}
\ExtensionTok{./tools/measure\_stack.sh} \AttributeTok{{-}{-}role}\NormalTok{ edge}
\end{Highlighting}
\end{Shaded}

Гаралтаас дараах хүснэгтийг бөглө:

\hypertarget{ux445ux4afux441ux43dux44dux433ux442-1.1-pi-3b-ux433ux438ux439ux43d-ux441ux443ux443ux440ux44c-ux442ux4e9ux43bux4e9ux432-ux44eux443-ux447-ux430ux436ux438ux43bux43bux430ux430ux433ux4afux439}{%
\subsubsection{Хүснэгт 1.1 --- Pi 3B-гийн суурь төлөв (юу ч ажиллаагүй)}\label{ux445ux4afux441ux43dux44dux433ux442-1.1-pi-3b-ux433ux438ux439ux43d-ux441ux443ux443ux440ux44c-ux442ux4e9ux43bux4e9ux432-ux44eux443-ux447-ux430ux436ux438ux43bux43bux430ux430ux433ux4afux439}}

\begingroup\scriptsize\begin{longtable}[]{@{}lll@{}}
\toprule
Хэмжигдэхүүн & Утга & Тэмдэглэл \\
\midrule
\endhead
\bottomrule
\endlastfoot
Pi загвар & & \texttt{Raspberry\ Pi\ 3\ Model\ B\ …} \\
\texttt{MemTotal} (MiB) & & \textasciitilde925 хүлээгдэнэ \\
\texttt{MemAvailable} (MiB) & & \\
Swap ашигласан (MiB) & & 0 байх ёстой \\
CPU температур (°C) & & \\
\texttt{get\_throttled} & & \textbf{0x0} байх ёстой \\
microSD чөлөөтэй (GB) & & \\
Ачаалал (1 мин) & & \\
\end{longtable}\endgroup

\begin{quote}
\textbf{\texttt{gpu\_mem} шалгалт:} \texttt{vcgencmd\ get\_mem\ gpu} → \texttt{16M} байх ёстой. \texttt{64M} бол SETUP.md Б.3 хийгээгүй. 48 MiB алдаж байна.
\end{quote}

\hypertarget{ux430ux43bux445ux430ux43c-2-ux4afux4afux43bux43dux438ux439-ux434ux430ux432ux445ux430ux440ux433ux44bux433-ux430ux441ux430ux430ux445-ux431ux430-ux445ux44dux43cux436ux438ux445-35-ux43cux438ux43d-ux437ux43a}{%
\subsection{Алхам 2 --- Үүлний давхаргыг асаах ба хэмжих (35 мин) {[}ЗК{]}}\label{ux430ux43bux445ux430ux43c-2-ux4afux4afux43bux43dux438ux439-ux434ux430ux432ux445ux430ux440ux433ux44bux433-ux430ux441ux430ux430ux445-ux431ux430-ux445ux44dux43cux436ux438ux445-35-ux43cux438ux43d-ux437ux43a}}

\begin{Shaded}
\begin{Highlighting}[]
\BuiltInTok{cd}\NormalTok{ stack}
\ExtensionTok{docker}\NormalTok{ compose }\AttributeTok{{-}{-}profile}\NormalTok{ core down }\DecValTok{2}\OperatorTok{\textgreater{}}\NormalTok{/dev/null}
\ExtensionTok{./../tools/measure\_stack.sh} \AttributeTok{{-}{-}role}\NormalTok{ cloud }\AttributeTok{{-}{-}startup}
\end{Highlighting}
\end{Shaded}

\texttt{-\/-startup} нь \texttt{docker\ compose\ -\/\allowbreak{}-profile\ core\ up\ -d} ажиллуулж, үйлчилгээ бүрийн \textbf{эрүүл болох хүртэлх} хугацааг хэмжинэ.

\hypertarget{ux445ux4afux441ux43dux44dux433ux442-1.2-ux4afux4afux43bux43dux438ux439-ux4afux439ux43bux447ux438ux43bux433ux44dux44dux43dux438ux439-ux44dux445ux43bux44dux445-ux445ux443ux433ux430ux446ux430ux430}{%
\subsubsection{Хүснэгт 1.2 --- Үүлний үйлчилгээний эхлэх хугацаа}\label{ux445ux4afux441ux43dux44dux433ux442-1.2-ux4afux4afux43bux43dux438ux439-ux4afux439ux43bux447ux438ux43bux433ux44dux44dux43dux438ux439-ux44dux445ux43bux44dux445-ux445ux443ux433ux430ux446ux430ux430}}

\begingroup\scriptsize\begin{longtable}[]{@{}lllll@{}}
\toprule
Үйлчилгээ & Эхэлсэн (сек) & Эрүүл болсон (сек) & RAM (MiB) & CPU \% \\
\midrule
\endhead
\bottomrule
\endlastfoot
emqx & & & & \\
influxdb & & & & \\
grafana & & & & \\
registry & & & & \\
\textbf{Нийт} & & & & \\
\end{longtable}\endgroup

Шалгах:

\begin{Shaded}
\begin{Highlighting}[]
\FunctionTok{make}\NormalTok{ health              }\CommentTok{\# бүх мөр 200}
\ExtensionTok{docker}\NormalTok{ compose ps}
\end{Highlighting}
\end{Shaded}

Хөтчөөр нээ:

\begingroup\scriptsize\begin{longtable}[]{@{}lll@{}}
\toprule
Үйлчилгээ & Хаяг & Нэвтрэх \\
\midrule
\endhead
\bottomrule
\endlastfoot
EMQX самбар & http://localhost:18083 & admin / \texttt{.env}-ийн нууц үг \\
Grafana & http://localhost:3000 & admin / \texttt{.env}-ийн нууц үг \\
InfluxDB & http://localhost:8181/health & --- \\
Бүртгэл & http://localhost:8090/health & --- \\
\end{longtable}\endgroup

\hypertarget{ux430ux43bux445ux430ux43c-3-ux438ux440ux43cux44dux433ux438ux439ux43d-ux434ux430ux432ux445ux430ux440ux433ux44bux433-ux430ux441ux430ux430ux445-25-ux43cux438ux43d-pi}{%
\subsection{Алхам 3 --- Ирмэгийн давхаргыг асаах (25 мин) {[}Pi{]}}\label{ux430ux43bux445ux430ux43c-3-ux438ux440ux43cux44dux433ux438ux439ux43d-ux434ux430ux432ux445ux430ux440ux433ux44bux433-ux430ux441ux430ux430ux445-25-ux43cux438ux43d-pi}}

\begin{Shaded}
\begin{Highlighting}[]
\BuiltInTok{cd}\NormalTok{ \textasciitilde{}/cnc302/edge}
\FunctionTok{cat}\NormalTok{ .env }\KeywordTok{|} \FunctionTok{grep}\NormalTok{ CLOUD\_HOST       }\CommentTok{\# компьютерийн LAN IP мөн үү?}
\FunctionTok{make}\NormalTok{ bridge}
\FunctionTok{make}\NormalTok{ up}
\end{Highlighting}
\end{Shaded}

Гүүр холбогдсоныг \textbf{хоёр талаас} батал:

\begin{Shaded}
\begin{Highlighting}[]
\CommentTok{\# [Pi] Pi дээр}
\FunctionTok{make}\NormalTok{ link}
\CommentTok{\# → cnc302/shutis/mhts/lab/pi3b{-}01/bridge/state 1}
\end{Highlighting}
\end{Shaded}

\begin{Shaded}
\begin{Highlighting}[]
\CommentTok{\# [ЗК] компьютер дээр}
\ExtensionTok{mosquitto\_sub} \AttributeTok{{-}h}\NormalTok{ localhost }\AttributeTok{{-}t} \StringTok{\textquotesingle{}cnc302/\#\textquotesingle{}} \AttributeTok{{-}v}
\end{Highlighting}
\end{Shaded}

\texttt{1} гарахгүй бол \texttt{docs/troubleshooting.md} §1.

Одоо ирмэгийн агентыг ажиллуул ({[}Pi{]} өөр терминал):

\begin{Shaded}
\begin{Highlighting}[]
\BuiltInTok{cd}\NormalTok{ \textasciitilde{}/cnc302/edge }\KeywordTok{\&\&} \FunctionTok{make}\NormalTok{ agent}
\end{Highlighting}
\end{Shaded}

Компьютер дээрх \texttt{mosquitto\_sub} цонхонд телеметр урсаж эхлэх ёстой. \textbf{Хэрэв урсахгүй бол §2 троублшүүтинг --- сэдвийн угтвар.}

\hypertarget{ux445ux4afux441ux43dux44dux433ux442-1.3-ux438ux440ux43cux44dux433ux438ux439ux43d-ux4afux439ux43bux447ux438ux43bux433ux44dux44dux43dux438ux439-ux43dux4e9ux4e9ux446}{%
\subsubsection{Хүснэгт 1.3 --- Ирмэгийн үйлчилгээний нөөц}\label{ux445ux4afux441ux43dux44dux433ux442-1.3-ux438ux440ux43cux44dux433ux438ux439ux43d-ux4afux439ux43bux447ux438ux43bux433ux44dux44dux43dux438ux439-ux43dux4e9ux4e9ux446}}

\begingroup\scriptsize\begin{longtable}[]{@{}llll@{}}
\toprule
Хэсэг & RAM (MiB) & CPU \% & Тэмдэглэл \\
\midrule
\endhead
\bottomrule
\endlastfoot
mosquitto & & & \\
ирмэгийн агент & & & \\
Docker демон & & & \\
\textbf{Нийт нэмэгдсэн} & & & Хүснэгт 1.1-тэй харьцуул \\
\textbf{Үлдэгдэл \texttt{MemAvailable}} & & & \\
\end{longtable}\endgroup

\hypertarget{ux430ux43bux445ux430ux43c-4-ux433ux4afux4afux440ux43dux438ux439-ux441ux430ux430ux442ux43bux44bux433-ux445ux44dux43cux436ux438ux445-40-ux43cux438ux43d-piux437ux43a}{%
\subsection{Алхам 4 --- Гүүрний саатлыг хэмжих (40 мин) {[}Pi{]}{[}ЗК{]}}\label{ux430ux43bux445ux430ux43c-4-ux433ux4afux4afux440ux43dux438ux439-ux441ux430ux430ux442ux43bux44bux433-ux445ux44dux43cux436ux438ux445-40-ux43cux438ux43d-piux437ux43a}}

Гурван замын саатлыг харьцуулна. Энэ бол Лаб 3-ын бүрэн хэмжилтийн \textbf{урьдчилсан хувилбар}.

\textbf{А. Loopback (сүлжээгүй)} --- {[}Pi{]} Pi дээр, Pi-гийн өөрийн брокер руу:

\begin{Shaded}
\begin{Highlighting}[]
\BuiltInTok{cd}\NormalTok{ \textasciitilde{}/cnc302}
\ExtensionTok{python3}\NormalTok{ tools/qos\_latency.py }\AttributeTok{{-}{-}host}\NormalTok{ localhost }\AttributeTok{{-}{-}port}\NormalTok{ 1883 }\DataTypeTok{\textbackslash{}}
        \AttributeTok{{-}{-}path}\NormalTok{ loopback }\AttributeTok{{-}{-}qos}\NormalTok{ 1 }\AttributeTok{{-}{-}count}\NormalTok{ 300 }\AttributeTok{{-}{-}interval}\NormalTok{ 0.05}
\end{Highlighting}
\end{Shaded}

\textbf{Б. LAN (шууд үүл рүү, гүүрийг тойрч)} --- {[}Pi{]} Pi дээр:

\begin{Shaded}
\begin{Highlighting}[]
\ExtensionTok{python3}\NormalTok{ tools/qos\_latency.py }\AttributeTok{{-}{-}host} \VariableTok{$CLOUD\_HOST} \AttributeTok{{-}{-}port}\NormalTok{ 1883 }\DataTypeTok{\textbackslash{}}
        \AttributeTok{{-}{-}path}\NormalTok{ lan }\AttributeTok{{-}{-}qos}\NormalTok{ 1 }\AttributeTok{{-}{-}count}\NormalTok{ 300 }\AttributeTok{{-}{-}interval}\NormalTok{ 0.05}
\end{Highlighting}
\end{Shaded}

\textbf{В. Гүүрээр} --- {[}Pi{]} нийтэлж, {[}ЗК{]} хүлээн авна. Хамгийн энгийн хэмжих арга:

\begin{Shaded}
\begin{Highlighting}[]
\CommentTok{\# [Pi] Pi дээр — цуурай хийх (echo) хэрэгсэл}
\ExtensionTok{mosquitto\_sub} \AttributeTok{{-}h}\NormalTok{ localhost }\AttributeTok{{-}t} \StringTok{\textquotesingle{}cnc302/shutis/mhts/lab/pi3b{-}01/cmd\textquotesingle{}} \AttributeTok{{-}v} \KeywordTok{\&}
\CommentTok{\# [ЗК] компьютер дээр — цагийн тэмдэгтэй мессеж илгээж, буцаж ирэхийг хүлээх}
\ExtensionTok{python3}\NormalTok{ tools/qos\_latency.py }\AttributeTok{{-}{-}host}\NormalTok{ localhost }\AttributeTok{{-}{-}port}\NormalTok{ 1883 }\DataTypeTok{\textbackslash{}}
        \AttributeTok{{-}{-}path}\NormalTok{ bridge }\AttributeTok{{-}{-}qos}\NormalTok{ 1 }\AttributeTok{{-}{-}count}\NormalTok{ 300 }\AttributeTok{{-}{-}interval}\NormalTok{ 0.05}
\end{Highlighting}
\end{Shaded}

\hypertarget{ux445ux4afux441ux43dux44dux433ux442-1.4-ux433ux443ux440ux432ux430ux43d-ux437ux430ux43cux44bux43d-ux441ux430ux430ux442ux430ux43b-qos-1-300-ux43cux435ux441ux441ux435ux436}{%
\subsubsection{Хүснэгт 1.4 --- Гурван замын саатал (QoS 1, 300 мессеж)}\label{ux445ux4afux441ux43dux44dux433ux442-1.4-ux433ux443ux440ux432ux430ux43d-ux437ux430ux43cux44bux43d-ux441ux430ux430ux442ux430ux43b-qos-1-300-ux43cux435ux441ux441ux435ux436}}

\begingroup\scriptsize\begin{longtable}[]{@{}llllll@{}}
\toprule
Зам & p50 (мс) & p95 (мс) & \textbf{p99 (мс)} & max (мс) & Алдагдал \% \\
\midrule
\endhead
\bottomrule
\endlastfoot
loopback (Pi доторх) & & & & & \\
LAN (Pi → EMQX) & & & & & \\
гүүрээр & & & & & \\
\end{longtable}\endgroup

\begin{quote}
\textbf{p99 нь гол тоо, дундаж биш.} Pi 3B-гийн 100 Mbit холбоос USB 2.0-ийн зурвасыг хуваадаг тул p50 сайхан харагдаж байхад p99 хэд дахин том гарч болно. Хэрэглэгчийн мэдэрдэг зүйл нь p99.
\end{quote}

\hypertarget{ux430ux43bux445ux430ux43c-5-ux43dux4e9ux4e9ux446ux438ux439ux43d-ux445ux44fux437ux433ux430ux430ux440ux44bux433-ux438ux43b-ux433ux430ux440ux433ux430ux445-30-ux43cux438ux43d-pi}{%
\subsection{Алхам 5 --- Нөөцийн хязгаарыг ил гаргах (30 мин) {[}Pi{]}}\label{ux430ux43bux445ux430ux43c-5-ux43dux4e9ux4e9ux446ux438ux439ux43d-ux445ux44fux437ux433ux430ux430ux440ux44bux433-ux438ux43b-ux433ux430ux440ux433ux430ux445-30-ux43cux438ux43d-pi}}

Виртуал флотыг Pi-гийн брокер руу чиглүүлж, санах ой хаана ханахыг \textbf{бодитоор} үзнэ.

\begin{Shaded}
\begin{Highlighting}[]
\CommentTok{\# [ЗК] компьютер дээрээс Pi{-}гийн брокер руу флот илгээнэ}
\ExtensionTok{python3}\NormalTok{ tools/sim\_device.py }\AttributeTok{{-}{-}host} \OperatorTok{\textless{}}\NormalTok{PI{-}IP}\OperatorTok{\textgreater{}}\NormalTok{ {-}{-}port 1883 }\DataTypeTok{\textbackslash{}}
        \AttributeTok{{-}{-}target}\NormalTok{ edge }\AttributeTok{{-}{-}devices}\NormalTok{ 50 }\AttributeTok{{-}{-}interval}\NormalTok{ 1.0}

\CommentTok{\# [Pi] зэрэгцүүлэн хэмжинэ}
\ExtensionTok{./tools/measure\_stack.sh} \AttributeTok{{-}{-}role}\NormalTok{ edge }\AttributeTok{{-}{-}watch}\NormalTok{ 120 5}
\end{Highlighting}
\end{Shaded}

Төхөөрөмжийн тоог 50 → 100 → 200 болгож давтана. Аль үед \texttt{MemAvailable} 150 MiB-аас доош унах вэ?

\hypertarget{ux445ux4afux441ux43dux44dux433ux442-1.5-ux438ux440ux43cux44dux433ux438ux439ux43d-ux430ux447ux430ux430ux43bux43bux44bux43d-ux445ux430ux440ux438ux443-ux4afux439ux43bux434ux44dux43b}{%
\subsubsection{Хүснэгт 1.5 --- Ирмэгийн ачааллын хариу үйлдэл}\label{ux445ux4afux441ux43dux44dux433ux442-1.5-ux438ux440ux43cux44dux433ux438ux439ux43d-ux430ux447ux430ux430ux43bux43bux44bux43d-ux445ux430ux440ux438ux443-ux4afux439ux43bux434ux44dux43b}}

\begingroup\scriptsize\begin{longtable}[]{@{}
  >{\raggedright\arraybackslash}p{(\columnwidth - 10\tabcolsep) * \real{0.1667}}
  >{\raggedright\arraybackslash}p{(\columnwidth - 10\tabcolsep) * \real{0.1667}}
  >{\raggedright\arraybackslash}p{(\columnwidth - 10\tabcolsep) * \real{0.1667}}
  >{\raggedright\arraybackslash}p{(\columnwidth - 10\tabcolsep) * \real{0.1667}}
  >{\raggedright\arraybackslash}p{(\columnwidth - 10\tabcolsep) * \real{0.1667}}
  >{\raggedright\arraybackslash}p{(\columnwidth - 10\tabcolsep) * \real{0.1667}}@{}}
\toprule
\begin{minipage}[b]{\linewidth}\raggedright
Виртуал төхөөрөмж
\end{minipage} & \begin{minipage}[b]{\linewidth}\raggedright
mosquitto RAM (MiB)
\end{minipage} & \begin{minipage}[b]{\linewidth}\raggedright
\texttt{MemAvailable} (MiB)
\end{minipage} & \begin{minipage}[b]{\linewidth}\raggedright
Swap \texttt{si/so}
\end{minipage} & \begin{minipage}[b]{\linewidth}\raggedright
Темп (°C)
\end{minipage} & \begin{minipage}[b]{\linewidth}\raggedright
\texttt{get\_throttled}
\end{minipage} \\
\midrule
\endhead
\bottomrule
\endlastfoot
0 (суурь) & & & & & \\
50 & & & & & \\
100 & & & & & \\
200 & & & & & \\
\end{longtable}\endgroup

\begin{quote}
⚠ Swap идэвхжвэл (\texttt{si/so} \textgreater{} 0) \textbf{зогсоо}. Тэр бол хязгаар. Цааш шахах нь Pi-г царцаана. Энэ нь алдаа биш, \textbf{хэмжилт}.
\end{quote}

\hypertarget{ux430ux43bux445ux430ux43c-6-ux44fux430ux433ux430ux430ux434-ux43fux43bux430ux442ux444ux43eux440ux43cux44bux433-pi-ux434ux44dux44dux440-ux430ux436ux438ux43bux43bux443ux443ux43bux436-ux431ux43eux43bux43eux445ux433ux4afux439-ux432ux44d-20-ux43cux438ux43d}{%
\subsection{Алхам 6 --- Яагаад платформыг Pi дээр ажиллуулж болохгүй вэ (20 мин)}\label{ux430ux43bux445ux430ux43c-6-ux44fux430ux433ux430ux430ux434-ux43fux43bux430ux442ux444ux43eux440ux43cux44bux433-pi-ux434ux44dux44dux440-ux430ux436ux438ux43bux43bux443ux443ux43bux436-ux431ux43eux43bux43eux445ux433ux4afux439-ux432ux44d-20-ux43cux438ux43d}}

Таамаглах биш, \textbf{тоолох}. {[}ЗК{]} компьютер дээр:

\begin{Shaded}
\begin{Highlighting}[]
\ExtensionTok{docker}\NormalTok{ stats }\AttributeTok{{-}{-}no{-}stream} \AttributeTok{{-}{-}format} \StringTok{\textquotesingle{}table \{\{.Name\}\}\textbackslash{}t\{\{.MemUsage\}\}\textquotesingle{}}
\end{Highlighting}
\end{Shaded}

\hypertarget{ux445ux4afux441ux43dux44dux433ux442-1.6-ux445ux44dux440ux44dux432-ux431ux4afux433ux434ux438ux439ux433-pi-ux434ux44dux44dux440-ux430ux436ux438ux43bux43bux443ux443ux43bux430ux445-ux433ux44dux432ux44dux43b}{%
\subsubsection{Хүснэгт 1.6 --- Хэрэв бүгдийг Pi дээр ажиллуулах гэвэл}\label{ux445ux4afux441ux43dux44dux433ux442-1.6-ux445ux44dux440ux44dux432-ux431ux4afux433ux434ux438ux439ux433-pi-ux434ux44dux44dux440-ux430ux436ux438ux43bux43bux443ux443ux43bux430ux445-ux433ux44dux432ux44dux43b}}

\begingroup\scriptsize\begin{longtable}[]{@{}lll@{}}
\toprule
Үйлчилгээ & Компьютер дээрх бодит RAM (MiB) & Pi-д багтах уу \\
\midrule
\endhead
\bottomrule
\endlastfoot
emqx & & \\
influxdb & & \\
grafana & & \\
registry & & \\
\textbf{Нийт} & & \\
Pi-д боломжтой (Хүснэгт 1.1) & & \\
\textbf{Дутагдал} & & \\
\end{longtable}\endgroup

Нэмэлт (сонголтот, {[}ЗК{]} хангалттай RAM байвал): ThingsBoard-ыг асааж хэмж.

\begin{Shaded}
\begin{Highlighting}[]
\ExtensionTok{docker}\NormalTok{ compose }\AttributeTok{{-}f}\NormalTok{ docker{-}compose.yml }\AttributeTok{{-}f}\NormalTok{ docker{-}compose.tb.yml }\DataTypeTok{\textbackslash{}}
  \AttributeTok{{-}{-}profile}\NormalTok{ core }\AttributeTok{{-}{-}profile}\NormalTok{ tb up }\AttributeTok{{-}d}
\CommentTok{\# 2{-}3 минут хүлээгээд:}
\ExtensionTok{docker}\NormalTok{ stats }\AttributeTok{{-}{-}no{-}stream}\NormalTok{ cnc302{-}thingsboard}
\end{Highlighting}
\end{Shaded}

Энэ нэг тоо (1.2--1.8 GiB) нь Pi-гийн \textbf{нийт} санах ойноос хэдэн дахин их вэ?

\hypertarget{ux430ux43bux445ux430ux43c-7-grafana-ux434-ux430ux43cux44cux434-ux441ux430ux43cux431ux430ux440-25-ux43cux438ux43d-ux437ux43a}{%
\subsection{Алхам 7 --- Grafana-д амьд самбар (25 мин) {[}ЗК{]}}\label{ux430ux43bux445ux430ux43c-7-grafana-ux434-ux430ux43cux44cux434-ux441ux430ux43cux431ux430ux440-25-ux43cux438ux43d-ux437ux43a}}

\begin{enumerate}
\def\labelenumi{\arabic{enumi}.}
\tightlist
\item
  http://localhost:3000 → Connections → Data sources → InfluxDB аль хэдийн бүртгэгдсэн (provisioning)
\item
  New dashboard → 3 панел:

  \begin{itemize}
  \tightlist
  \item
    Pi-гийн CPU температур (\texttt{health} сувгаас)
  \item
    \texttt{MemAvailable} (\texttt{health} сувгаас)
  \item
    Гүүрний төлөв (\texttt{bridge/state}) --- Stat панел
  \end{itemize}
\item
  Хадгалж, JSON-ыг экспортлон \texttt{lab01/dashboard.json} болгож commit хийнэ
\end{enumerate}

\begin{quote}
Энэ самбар Лаб 5-д гүүр тасрахыг \textbf{нүдээр харах} гол хэрэгсэл болно.
\end{quote}

\hypertarget{ux430ux43bux445ux430ux43c-8-git-commit-10-ux43cux438ux43d}{%
\subsection{Алхам 8 --- Git commit (10 мин)}\label{ux430ux43bux445ux430ux43c-8-git-commit-10-ux43cux438ux43d}}

\begin{Shaded}
\begin{Highlighting}[]
\BuiltInTok{cd}\NormalTok{ \textasciitilde{}/cnc302}
\FunctionTok{git}\NormalTok{ add stack/ edge/ lab01/}
\FunctionTok{git}\NormalTok{ commit }\AttributeTok{{-}m} \StringTok{"Лаб 1: хоёр давхаргат платформ, гүүр ажиллаж байна"}
\FunctionTok{git}\NormalTok{ tag lab01{-}done}
\FunctionTok{git}\NormalTok{ push }\KeywordTok{\&\&} \FunctionTok{git}\NormalTok{ push }\AttributeTok{{-}{-}tags}
\end{Highlighting}
\end{Shaded}

⚠ \texttt{.env}, \texttt{bridge.conf} \textbf{орохгүй} байгааг шалга:

\begin{Shaded}
\begin{Highlighting}[]
\FunctionTok{git}\NormalTok{ status }\AttributeTok{{-}{-}short}
\FunctionTok{git}\NormalTok{ ls{-}files }\KeywordTok{|} \FunctionTok{grep} \AttributeTok{{-}E} \StringTok{\textquotesingle{}\textbackslash{}.env$|bridge\textbackslash{}.conf$\textquotesingle{}}    \CommentTok{\# хоосон байх ЁСТОЙ}
\end{Highlighting}
\end{Shaded}

\hypertarget{ux445ux44fux43dux430ux43bux442ux44bux43d-ux430ux441ux443ux443ux43bux442ux443ux443ux434}{%
\section{5. Хяналтын асуултууд}\label{ux445ux44fux43dux430ux43bux442ux44bux43d-ux430ux441ux443ux443ux43bux442ux443ux443ux434}}

Тайландаа \textbf{тоон баримт заавал иш татаж} хариулна.

\begin{enumerate}
\def\labelenumi{\arabic{enumi}.}
\item
  Хүснэгт 1.4-д гурван замын p99 хэд дахин ялгаатай вэ? Ялгааны \textbf{гол шалтгаан} юу вэ --- сүлжээ, брокерын боловсруулалт, эсвэл дараалал? Хэрхэн ялгаж тогтоох вэ?
\item
  Хүснэгт 1.5-д хязгаарлагч нөөц юу байв --- RAM, CPU, эсвэл сүлжээ? Хариултаа тоогоор нотол. (Санамж: \texttt{docker\ stats}-ын CPU \% нь 4 цөм дээр 400\% хүртэл гарна.)
\item
  Хүснэгт 1.6-гийн дутагдал хэдэн MiB вэ? Хэрэв Pi 3B-д 2 GB RAM байсан бол платформыг түүн дээр ажиллуулах уу? Яагаад тийм / үгүй?
\item
  \texttt{cleansession\ false}-ийг \texttt{true} болговол юу өөрчлөгдөх вэ? Ямар тохиолдолд \texttt{true} нь \textbf{зөв} сонголт байх вэ?
\item
  Гүүр нь \texttt{cnc302/shutis/\#} сэдвийг дамжуулдаг. Хэрэв өөр баг \texttt{SITE=shutis} ашиглаад нэг брокерт холбогдвол юу болох вэ? Үүнээс хэрхэн сэргийлэх вэ (хоёр өөр арга нэрлэ)?
\item
  Pi дээр \texttt{vcgencmd\ get\_throttled} нь \texttt{0x50000} буцаав. Энэ юу гэсэн үг вэ, өмнөх хэмжилтүүд хүчинтэй юу?
\end{enumerate}

\hypertarget{ux445ux4afux43bux44dux44dux43bux433ux44dux43d-ux4e9ux433ux4e9ux445-ux437ux4afux439ux43b}{%
\section{6. Хүлээлгэн өгөх зүйл}\label{ux445ux4afux43bux44dux44dux43bux433ux44dux43d-ux4e9ux433ux4e9ux445-ux437ux4afux439ux43b}}

\begingroup\scriptsize\begin{longtable}[]{@{}lll@{}}
\toprule
\# & Зүйл & Байрлал \\
\midrule
\endhead
\bottomrule
\endlastfoot
1 & Тайлан (Хүснэгт 1.1--1.6 бөглөсөн) & \texttt{lab01/report.md} → PDF \\
2 & Grafana самбарын JSON & \texttt{lab01/dashboard.json} \\
3 & \texttt{measure\_stack.sh}-ийн түүхий гаралт & \texttt{lab01/out/} \\
4 & Git tag & \texttt{lab01-done} \\
\end{longtable}\endgroup

\hypertarget{ux4afux43dux44dux43bux433ux44dux44dux43dux438ux439-ux448ux430ux43bux433ux443ux443ux440-10-ux43eux43dux43eux43e}{%
\section{7. Үнэлгээний шалгуур (10 оноо)}\label{ux4afux43dux44dux43bux433ux44dux44dux43dux438ux439-ux448ux430ux43bux433ux443ux443ux440-10-ux43eux43dux43eux43e}}

\begingroup\scriptsize\begin{longtable}[]{@{}ll@{}}
\toprule
Шалгуур & Оноо \\
\midrule
\endhead
\bottomrule
\endlastfoot
Хоёр давхарга ажиллаж, гүүр холбогдсон (амьд үзүүлнэ) & 3 \\
Хүснэгт 1.1--1.6 бүрэн, throttling шалгасан & 3 \\
Хязгаарлагч нөөцийг \textbf{зөв тогтоож, тоогоор нотолсон} & 2 \\
Git сан цэвэр, нууц файл ороогүй & 1 \\
Grafana самбар ажиллаж байна & 1 \\
\end{longtable}\endgroup

\hypertarget{ux442ux4afux433ux44dux44dux43cux44dux43b-ux430ux43bux434ux430ux430}{%
\section{8. Түгээмэл алдаа}\label{ux442ux4afux433ux44dux44dux43cux44dux43b-ux430ux43bux434ux430ux430}}

\begingroup\scriptsize\begin{longtable}[]{@{}
  >{\raggedright\arraybackslash}p{(\columnwidth - 4\tabcolsep) * \real{0.3333}}
  >{\raggedright\arraybackslash}p{(\columnwidth - 4\tabcolsep) * \real{0.3333}}
  >{\raggedright\arraybackslash}p{(\columnwidth - 4\tabcolsep) * \real{0.3333}}@{}}
\toprule
\begin{minipage}[b]{\linewidth}\raggedright
Шинж
\end{minipage} & \begin{minipage}[b]{\linewidth}\raggedright
Шалтгаан
\end{minipage} & \begin{minipage}[b]{\linewidth}\raggedright
Шийдэл
\end{minipage} \\
\midrule
\endhead
\bottomrule
\endlastfoot
\texttt{bridge/state\ 0} & CLOUD\_HOST буруу / галт хана & troubleshooting §1 \\
Гүүр холбогдсон ч мессеж алга & \texttt{SITE} зөрсөн & troubleshooting §2 \\
Санах ойн тоо хачин & VS Code сервер ажиллаж байна & \texttt{pkill\ -f\ vscode-server} \\
\texttt{docker\ stats} санах ой хоосон & cgroup идэвхгүй & SETUP.md Б.5 \\
Саатал 10 дахин хэлбэлзэнэ & throttling эсвэл Wi-Fi & \texttt{get\_throttled}, кабельд шилжих \\
\texttt{exec\ format\ error} & 32-bit OS & \texttt{uname\ -m} → aarch64 байх ёстой \\
\end{longtable}\endgroup
